\section{A classical analogue: staged objecthood in metastable Markov dynamics}\label{sec:markov}

SBT's central mechanism---that objecthood is layer-relative and stabilized by packaging---is not uniquely quantum.
A purely classical substrate with coarse access can exhibit the same structure: at certain timescales, coarse variables behave like stable objects; at other timescales, they do not.

\subsection{Metastable substrate and a basin lens}

Let $Z$ be a finite set of microstates, and let $P$ be a Markov kernel on~$Z$.\cite{Norris1997}
Write $T_\tau$ for $\tau$ steps of evolution acting on a microdistribution~$\mu$:
\[
T_\tau(\mu) = \mu P^\tau.
\]
Assume that $P$ has (at least) two metastable basins with weak coupling: rapid mixing within each basin and rare transitions between basins.

Define a basin lens (coarse record variable)
\[
f : Z \to X,
\]
where $X$ is a small set of macro labels (e.g., ``left basin'' versus ``right basin'').

\subsection{Prototypes and the packaging operator $E_{\tau,f}$}

Coarse access pushes a microdistribution forward to a distribution over basins:
\[
Q_f(\mu)(x) := \sum_{z\in Z: f(z)=x} \mu(z).
\]
Completion returns a canonical representative on~$Z$ by selecting a prototype distribution~$u_x$ for each macro label $x\in X$.
In the metastable example, prototypes are chosen as uniform distributions on each basin.
These prototypes (and hence~$U_f$) are a chosen completion/section and are part of the package, not canonical.

Define completion by
\[
U_f(p) := \sum_{x\in X} p(x)\,u_x,
\]
and define the timescale-dependent packaging operator by
\[
E_{\tau,f} := U_f \circ Q_f \circ T_\tau.
\]
That is, the operator evolves the substrate, reads only the basin masses, and then discards within-basin microdetails by re-imposing the prototypes.

\subsection{Idempotence defect and prototype stability}

Packaging behaves like a closure when applying it twice adds no new record-level content.
A quantitative way to test this property is via the idempotence defect:
\[
\delta(\tau) := \sup_{\mu}\ \mathrm{TV}\!\left(E_{\tau,f}(E_{\tau,f}(\mu)),\ E_{\tau,f}(\mu)\right),
\]
where $\mathrm{TV}$ denotes total variation distance on~$Z$.
When $\delta(\tau)$ is small, the packaged description is approximately closed under reapplication: the layer has stabilized a coherent object language at that timescale.

A complementary diagnostic is prototype stability: whether $E_{\tau,f}(u_x)\approx u_x$ for each macro label~$x$.
In our implementation, we track a worst-case stability score over prototypes as a function of~$\tau$.

\subsection{Two regimes: emergent objects and collapse-to-constant}

Figure~\ref{fig:markov-idem} plots the measured defect~$\delta(\tau)$ for a fixed metastable chain (repository experiment \texttt{EXP-MK1}; see Appendix~\ref{app:repro}).
The curve exhibits distinct regimes, reflecting the presence of multiple timescales:
\begin{itemize}
  \item At timescales where within-basin mixing dominates and cross-basin leakage is negligible, the basin prototypes behave as stable objects and packaging is close to idempotent.
  \item At intermediate timescales where leakage becomes significant, repeated packaging can drift: the basin split is not fully stable as an object language, and the defect increases.
  \item At sufficiently long timescales, the chain approaches global equilibrium, and packaging approaches a collapse-to-constant regime: the effective object language becomes trivial (essentially one stable macro-description), and the defect decreases again.
\end{itemize}
In the particular instance plotted, the first timescale meeting the chosen stability threshold occurs at $\tau=\MKFirstStableTau$ (see the experiment output).

We also ran robustness sweeps over multiple seeds for the full experimental pipeline; the Markov staging behavior is stable under these variations. A compact min/median/max summary is reported in Appendix~\ref{app:robustness}.

\begin{figure}[t]
  \centering
  \includegraphics[width=0.95\linewidth]{markov_idempotence_vs_tau.png}
  \caption{Idempotence defect $\delta(\tau)$ for a metastable Markov substrate under basin packaging. Distinct regimes reflect competing mixing timescales: emergent basin-level objects at some $\tau$, and a long-time collapse-to-constant regime.}
  \label{fig:markov-idem}
\end{figure}
